\chapter{Diskussion und Ausblick}\label{ch:discussion}

In diesem Kapitel werden die experimentellen Befunde vor dem Hintergrund der theoretischen Grundlagen und des Stands der Technik kritisch diskutiert. Die Analyse stützt sich auf die Verknüpfung der quantitativen Trainingsmetriken (Return, Episodenlänge und Entropieverlauf) über die fünf Zufalls-Seeds (1, 8, 42, 24 und 77) mit der qualitativen Verifikation der Bewegungsmuster anhand der aufgezeichneten Episoden. Darauf aufbauend werden die inhärenten Limitationen der Arbeit systematisch dargelegt und konkrete Ansatzpunkte für zukünftige Forschungsarbeiten abgeleitet.

\section{Zusammenfassung und Synthese der Ergebnisse}

\subsection{Interaktionsumgebung (Petting)}
Die Ergebnisse in der Interaktionsumgebung zeigen eine klare Zweiteilung:
\begin{itemize}
	\item \textbf{Erfolgreiche Basiskomponenten:} In der ersten Stufe erlernte die Policy das stabile Stehen zuverlässig über alle fünf Seeds ($N=5$), was sich in einer raschen Konvergenz der mittleren Episodenlänge auf das Maximum von $1000$ Zeitschritten und einem stabilen Haltungsreward widerspiegelt (vgl. \autoref{fig:stand_pet_metrics}). Die spätere Interaktionslogik nutzt ein positionsbasiertes Berührungssignal und zeigte in den aufgezeichneten Episoden eine Rückkehr in die neutrale Standpose nach Beendigung der Berührung.
	\item \textbf{Nicht erfolgreiche Gestenkoordination:} Das Erlernen einer zeitlich koordinierten, dynamischen Reaktionsgeste über die gesamte Dauer der Berührung war hingegen nicht erfolgreich. Trotz eines initialen Anstiegs der Episodenlänge fiel diese im weiteren Trainingsverlauf bei allen Seeds drastisch ab. Der Reward für die eigentliche Gestenbewegung ($r_{\text{gesture\_movement}}$) verharrte durchgehend nahe null, während die Policy-Entropie anstieg. Die Beobachtungen sind mit einem lokalen Optimum vereinbar, in dem ein kurzer Bewegungsimpuls gegenüber einer vollständigen Geste bevorzugt wird.
\end{itemize}

\subsection{Navigationsumgebung}
Das Curriculum-Training in der Navigationsumgebung demonstriert die Stärken und Grenzen des DRL-Paradigmas:
\begin{itemize}
	\item \textbf{Robuste Basislokomotion (Stufen 1 und 2):} Das Stehen und gerichtete Vorwärtslaufen wurden über alle Seeds hinweg mit hoher Stabilität und minimaler Varianz erlernt (\autoref{fig:stand_metrics} und \autoref{fig:walk_metrics}). Der steile Anstieg von Return und Episodenlänge um Iteration 300 lässt sich plausibel dadurch erklären, dass die Policy nach und nach die koordinierte Beinphasenverschiebung (diagonales Trabgangmuster) erlernt und dadurch Stolperer und Stürze vermieden werden, sodass die Episoden bis zur maximalen Schrittzahl durchlaufen und der kumulierte Return ansteigt. Die Konvergenz der Geschwindigkeits- und Stabilitätsterme ist mit einer verlässlichen Gangmustergenerierung vereinbar.
	\item \textbf{Passive Robustheit bei Hinderniskontakt (Stufe 3):} Bei der Konfrontation mit statischen Hindernissen zeigte die Policy eine hohe Haltungsstabilität. Kam es zu physischem Kontakt mit dem Hindernis, war die Laufpolicy ausreichend stabil, sodass der Roboter trotz des Kontakts mit dem Hindernis seine Körperhaltung beibehalten konnte. Das Vorbeigleiten am Hinderniskörper entstand als emergenter Nebeneffekt der vorwärtsgerichteten Geschwindigkeits- und Stabilitätsbelohnungen, nicht als vorausschauende perzeptive Pfadplanung.
	\item \textbf{Nicht erfolgreiche Sprungkoordination (Stufe 3b):} Das Erlernen explosiver Sprungmanöver über Hindernisse schlug fehl. Return und Episodenlänge oszillierten ohne Konvergenz, während die Entropie stark anstieg. Die Sprungbelohnungen waren zwar hoch skaliert, wurden aber nur unter engen räumlichen und dynamischen Bedingungen ausgelöst. Gleichzeitig lieferten die aktiven Lauf- und Stabilitätsbelohnungen häufiger verwertbares Feedback, sodass eine vollständige Sprungsequenz vermutlich zu selten als zusammenhängendes Lernsignal auftrat.
\end{itemize}

\section{Interpretation und Ursachenanalyse}

\subsection{Limitationen der Arbeit}\label{sec:limitations}
Die gewonnenen Erkenntnisse sind im Kontext folgender methodischer Limitationen zu interpretieren:
\begin{itemize}
	\item \textbf{Simulationsumgebung und Sim-to-Real-Transfer:} Sämtliche Experimente wurden ausschließlich in der Simulationsplattform Genesis durchgeführt. Ein physischer Transfer auf den realen Unitree Go2 Edu sowie der Einsatz von Domain Randomization (z.\,B. Variation von Reibungskoeffizienten, Massen und Latenzen) wurden im Rahmen dieser Arbeit nicht realisiert.
	\item \textbf{Statistischer Stichprobenumfang ($N=5$):} Die Verwendung von fünf Zufalls-Seeds ermöglicht die Identifikation grundlegender Konvergenz- und Divergenzmuster, reicht jedoch für strenge inferenzstatistische Signifikanzprüfungen (z.\,B. Hypothesentests mit engen Konfidenzintervallen) nicht aus.
	\item \textbf{Modellierung der Interaktionspartner:} In der Interaktionsumgebung wird die menschliche Hand als starre geometrische Kugel abstrahiert. Komplexe Nachgiebigkeiten, Reibungsphänomene und biomechanische Geometrien realer Hände werden hierbei nicht abgebildet.
	\item \textbf{Manuelle Curriculum-Transitionen:} Die Übergänge zwischen den Curriculum-Stufen erfolgten schrittweise manuell anhand empirischer Prüfkriterien anstelle automatisierter, metrikenbasierter Schwellenwertsteuerungen.
	\item \textbf{Hyperparameter- und Belohnungsabstimmung:} Die Netzwerkarchitektur mit drei Hidden Layern (512, 256 und 128 Neuronen) sowie die PPO-Hyperparameter wurden stufenabhängig festgelegt. Eine vollständige, automatisierte Bayes'sche Suche (z.\,B. via Optuna \cite{akibaOptunaNextgenerationHyperparameter2019}) wurde nicht für alle finalen Trainingsläufe durchgeführt.
\end{itemize}

\subsection{Ursachenanalyse der Interaktions- und Sprunginstabilitäten}

Die Ursachen unterscheiden sich zwischen Interaktions- und Sprungtraining. Bei der Interaktion wird der Kontakt nicht über eine Kraftmessung, sondern über den Abstand einer virtuellen Handkugel zum Kopf sowie zusätzliche Zustandsbedingungen erkannt. Das resultierende Signal ist binär und hängt von der zufälligen Kontaktverteilung ab. Während des Kontakts wird die Geste durch einen Timer aktiviert und beim Ende des Kontakts unmittelbar beendet. Dadurch entsteht kein klar abgegrenzter Lernabschnitt mit Start, Zwischenzielen und erfolgreichem Abschluss.

Die Divergenz im Spätverlauf der Standstufe (Stufe 1) lässt sich möglicherweise dadurch erklären, dass die Policy in der späten Explorationsphase stärkere Gelenkbewegungen testet, welche die kritischen Schwellenwerte für Roll- und Nickwinkel überschreiten und dadurch vorzeitige Episodenabbrüche auslösen. Dass die Seeds 1, 8 und 42 hierbei kollabieren, während die Seeds 24 und 77 stabil bleiben, deutet auf eine ausgeprägte seed-abhängige Trainingsinstabilität hin.

Zusätzlich belohnt $r_{\text{gesture\_movement}}$ den Betrag ausgewählter Gelenkgeschwindigkeiten, multipliziert mit einem aufrechten Standfaktor. Eine bestimmte Richtung, Amplitude oder Synchronisation der Geste wird nicht geprüft. Dem stehen starke Stabilitäts- und Haltungsbelohnungen gegenüber. Die Policy hat daher einen plausiblen Anreiz, die Bewegung klein zu halten, sobald zusätzliche Gelenkbewegungen das Sturzrisiko erhöhen. Der beobachtete Kollaps der Episodenlänge und die geringe Bewegungsbelohnung stützen diese Interpretation, beweisen aber keine einzelne kausale Ursache.

Beim Sprungtraining wird die Laufpolicy aus Stufe 2 mit einer wesentlich komplexeren Aufgabe konfrontiert. Der Roboter muss sich einem nur $2\,\text{m}$ entfernten Hindernis nähern, nahe am Hindernis einen Absprung mit $v_z>0{,}3\,\text{m/s}$ ausführen, zentriert bleiben, die erforderliche Rumpfhöhe erreichen und anschließend stabil landen. Die Sprungbelohnungen werden jeweils nur unter räumlich, höhen- und phasenabhängigen Bedingungen aktiviert. Ein leicht verspäteter, zu niedriger oder seitlich verschobener Versuch erzeugt deshalb nur wenig nutzbares Feedback.

Gleichzeitig bleiben mehrere Laufbelohnungen aktiv, während die seitliche Bewegung im Sprungcurriculum deaktiviert ist. Normales Vorwärtslaufen liefert damit häufiger und kontinuierlicher Feedback als ein vollständiger Sprung. Die oszillierenden Returns und die steigende Entropie sind folglich mit einer weiterhin hohen Aktionsunsicherheit vereinbar. Die beobachtete Verhaltenskonsequenz lässt sich daher nicht auf eine einzelne Reward-Komponente zurückführen: Obwohl die Sprungterme in Stufe 4 hoch skaliert sind, werden sie nur in einer engen räumlichen und dynamischen Trefferregion aktiviert, während die weiterhin aktiven Lauf- und Stabilitätsrewards kontinuierlich Feedback liefern. Das Verhalten der Policy war infolgedessen häufig defensiv, indem sie versuchte, das Hindernis frontal zu überklettern oder seitlich daran vorbeizulaufen.

\section{Praktische Relevanz und Ausblick}

Aus den Erkenntnissen dieser Arbeit ergeben sich mehrere vielversprechende Perspektiven für zukünftige Forschungsarbeiten:
\begin{itemize}
	\item \textbf{Sequenzielles Reward Shaping und Phasensteuerung:} Für die Gesten- und Sprungsteuerung sollten phasenabhängige Zustandsautomaten oder hierarchische RL-Architekturen implementiert werden, die Stabilitätsstrafen während der aktiven Flug- oder Gestenphase gezielt maskieren und erst bei der Landung reaktivieren.
	\item \textbf{Curriculum-Zwischenstufen:} Für dynamische Manöver wie das Springen empfiehlt sich die Einführung einer Zwischenstufe, in der der Roboter zunächst vertikale Absprünge und Landungen auf freiem, ebenem Boden ohne Hinderniskonfrontation erlernt, bevor das zeitliche Timing des Absprungs in Relation zum Hindernisabstand hinzukommt.
	\item \textbf{Imitation Learning und Demonstrationsdaten:} Die Kombination von RL mit Vorab-Demonstrationen (Behavior Cloning / Guided RL), beispielsweise aus Motion-Capture-Daten von Tieren oder analytischen Trajektorien (Bézier-Kurven), kann das Explorationsproblem bei komplexen Gesten und Sprüngen drastisch reduzieren.
	\item \textbf{Automatisierte Hyperparameteroptimierung und Reward-Design:} Der Einsatz automatisierter Optimierungsframeworks wie Optuna (mittels Tree-structured Parzen Estimator) oder LLM-gestützter Reward-Design-Pipelines wie Eureka bietet großes Potenzial, um die manuelle Belohnungsjustierung durch systematische Exploration zu ersetzen.
	\item \textbf{Sim-to-Real-Transfer:} Durch Integration umfassender Domain Randomization (Massen, Motorverzögerungen, Bodenreibung) sowie Sensorrauschen kann die Übertragbarkeit der in Genesis trainierten Basislauf- und Stabilitäts-Policies auf den physischen Unitree Go2 Edu in realen Laborumgebungen validiert werden.
\end{itemize}